\chapter{Operations on Sets}
\signature

\section{The basic operations}

Throughout, sets live inside a fixed universal set $U$.

\begin{definition}[Union, intersection, difference, complement]
\begin{align*}
A \cup B &= \{x \mid x \in A \vee x \in B\} && \text{(union)}\\
A \cap B &= \{x \mid x \in A \wedge x \in B\} && \text{(intersection)}\\
A \setminus B &= \{x \mid x \in A \wedge x \notin B\} && \text{(difference)}\\
\overline{A} &= U \setminus A = \{x \in U \mid x \notin A\} && \text{(complement)}\\
A \oplus B &= (A \setminus B) \cup (B \setminus A) && \text{(symmetric difference).}
\end{align*}
Two sets with $A \cap B = \varnothing$ are called \emph{disjoint}.
\end{definition}

Each set operation mirrors a logical connective: $\cup \leftrightarrow \vee$,
$\cap \leftrightarrow \wedge$, $\overline{\phantom{A}} \leftrightarrow \neg$,
$\oplus \leftrightarrow \oplus$. This dictionary transfers every law of
Chapter~2 into a set identity.

\begin{example}
With $U = \{1,\dots,10\}$, $A = \{1,2,3,4,5\}$, $B = \{4,5,6,7\}$:
$A \cup B = \{1,\dots,7\}$, $A \cap B = \{4,5\}$,
$A \setminus B = \{1,2,3\}$, $\overline{B} = \{1,2,3,8,9,10\}$,
$A \oplus B = \{1,2,3,6,7\}$.
\end{example}

\begin{theorem}[Inclusion--exclusion, basic form]
For finite sets, $|A \cup B| = |A| + |B| - |A \cap B|$.
\end{theorem}

\begin{example}[Combining operations]\label{ex:combo}
With the sets above, compute $\overline{A \cap B} \cap A$:
$A \cap B = \{4,5\}$, so $\overline{A \cap B} = U \setminus \{4,5\}$, and
intersecting with $A$ leaves $\{1,2,3\} = A \setminus B$. In fact
$\overline{A \cap B} \cap A = A \setminus B$ always --- provable by the
identities below.
\end{example}

\begin{example}[Empty regions]
$A \cap \overline{A} = \varnothing$ and $(A \setminus B) \cap (B \setminus A)
= \varnothing$ for any $A, B$: some Venn regions are empty by pure logic,
regardless of the particular sets.
\end{example}

\section{Set identities}

\begin{theorem}[Set identities]\label{thm:setids}
For all $A, B, C \subseteq U$:
\begin{center}
\renewcommand{\arraystretch}{1.3}
\begin{tabular}{ll}
\toprule
Identity & $A \cup \varnothing = A$,\quad $A \cap U = A$\\
Domination & $A \cup U = U$,\quad $A \cap \varnothing = \varnothing$\\
Idempotence & $A \cup A = A$,\quad $A \cap A = A$\\
Double complement & $\overline{\overline{A}} = A$\\
Commutativity & $A \cup B = B \cup A$,\quad $A \cap B = B \cap A$\\
Associativity & $(A \cup B) \cup C = A \cup (B \cup C)$, and dually\\
Distributivity & $A \cup (B \cap C) = (A \cup B) \cap (A \cup C)$, and dually\\
De Morgan & $\overline{A \cup B} = \overline{A} \cap \overline{B}$,\quad
            $\overline{A \cap B} = \overline{A} \cup \overline{B}$\\
Absorption & $A \cup (A \cap B) = A$,\quad $A \cap (A \cup B) = A$\\
Complement & $A \cup \overline{A} = U$,\quad $A \cap \overline{A} = \varnothing$\\
\bottomrule
\end{tabular}
\end{center}
Also $A \setminus B = A \cap \overline{B}$, the workhorse for eliminating
differences.
\end{theorem}

\section{Proving identities}

Three standard methods.

\textbf{1. The element method (double inclusion).} To prove $X = Y$, show
$x \in X \Rightarrow x \in Y$ and conversely, unfolding the definitions into
logic and applying the propositional laws.

\begin{example}
Prove $\overline{A \cup B} = \overline{A} \cap \overline{B}$:
\[
x \in \overline{A \cup B}
\iff \neg(x \in A \vee x \in B)
\iff (x \notin A) \wedge (x \notin B)
\iff x \in \overline{A} \cap \overline{B},
\]
the middle step being De Morgan's law for propositions.
\end{example}

\textbf{2. Algebraic manipulation.} Chain identities from
Theorem~\ref{thm:setids}, replacing equals by equals --- as in
Example~\ref{ex:combo}:
$\overline{A \cap B} \cap A
= (\overline{A} \cup \overline{B}) \cap A
= (\overline{A} \cap A) \cup (\overline{B} \cap A)
= \varnothing \cup (A \cap \overline{B}) = A \setminus B$.

\textbf{3. Membership tables.} Analogous to truth tables: each row records
whether $x$ belongs ($1$) or not ($0$) to each set; two expressions are equal
iff their columns coincide.

\begin{example}[Membership-table proof]
$A \oplus B = (A \cup B) \setminus (A \cap B)$:
\begin{center}
\begin{tabular}{cc|ccccc}
\toprule
$A$&$B$&$A\cup B$&$A\cap B$&$(A\cup B)\setminus(A\cap B)$&$A\oplus B$\\
\midrule
1&1&1&1&0&0\\
1&0&1&0&1&1\\
0&1&1&0&1&1\\
0&0&0&0&0&0\\
\bottomrule
\end{tabular}
\end{center}
The last two columns agree, proving the identity.
\end{example}

\begin{notebox}
A membership table over $n$ set letters has $2^n$ rows --- it is exactly a
truth table in disguise, so the method is complete for identities built from
$\cup, \cap, \setminus, \oplus$ and complement.
\end{notebox}

\section{Generalised unions and intersections}

For a family $A_1, \ldots, A_n$ (or an infinite family):
\[
\bigcup_{i=1}^{n} A_i = \{x \mid \exists i,\ x \in A_i\},
\qquad
\bigcap_{i=1}^{n} A_i = \{x \mid \forall i,\ x \in A_i\}.
\]

\begin{example}
With $A_i = \{i, i+1, i+2, \ldots\} \subseteq \N$:
$\bigcup_{i \geq 1} A_i = A_1$ and $\bigcap_{i \geq 1} A_i = \varnothing$
(no natural number belongs to every $A_i$).
\end{example}

\section{Summary}

Union, intersection, difference, complement and symmetric difference mirror
$\vee, \wedge$, $\setminus$, $\neg, \oplus$; consequently the identity,
domination, idempotent, distributive, De Morgan, absorption and complement
laws all transfer from logic to sets. Identities are proved by the element
method, by algebraic chaining, or mechanically by membership tables.

\section*{Exercises}
\addcontentsline{toc}{section}{Exercises}

\begin{exercise}{\easy}
Let $U = \{1,\dots,12\}$, $A = \{1,3,5,7,9,11\}$, $B = \{3,6,9,12\}$.
Compute $A \cup B$, $A \cap B$, $A \setminus B$, $B \setminus A$,
$\overline{A}$ and $A \oplus B$.
\end{exercise}

\begin{exercise}{\easy}
Draw (or describe region by region) the Venn diagram of
$(A \cup B) \setminus (A \cap B)$ and name the operation it represents.
\end{exercise}

\begin{exercise}{\easy}
Two clubs have $28$ and $19$ members; $7$ people belong to both. How many
people belong to at least one club?
\end{exercise}

\begin{exercise}{\medium}
Prove by the element method: $A \setminus (B \cup C) =
(A \setminus B) \cap (A \setminus C)$.
\end{exercise}

\begin{exercise}{\medium}
Prove algebraically (citing identities): $A \cup (B \setminus A) = A \cup B$.
\end{exercise}

\begin{exercise}{\medium}
Verify by a membership table:
$A \setminus (A \setminus B) = A \cap B$.
\end{exercise}

\begin{exercise}{\medium}
Show that $A \subseteq B$ if and only if $A \cup B = B$, and if and only if
$A \cap B = A$.
\end{exercise}

\begin{exercise}{\hard}
Prove that symmetric difference is associative:
$(A \oplus B) \oplus C = A \oplus (B \oplus C)$.
\emph{Hint: a membership table over three letters, or the parity
characterisation.}
\end{exercise}

\begin{exercise}{\hard}
Solve for $X$: the equation $A \oplus X = B$ has exactly one solution for
given $A, B \subseteq U$. Find it and prove uniqueness.
\end{exercise}

\begin{exercise}{\hard}
Prove the three-set inclusion--exclusion formula:
\[
|A \cup B \cup C| = |A|+|B|+|C| - |A\cap B| - |A\cap C| - |B\cap C|
+ |A\cap B\cap C|.
\]
\end{exercise}

\section*{Solutions to the Exercises}
\addcontentsline{toc}{section}{Solutions to the Exercises}

\begin{solution}{7.1}
$A \cup B = \{1,3,5,6,7,9,11,12\}$;\
$A \cap B = \{3,9\}$;\
$A \setminus B = \{1,5,7,11\}$;\
$B \setminus A = \{6,12\}$;\
$\overline{A} = \{2,4,6,8,10,12\}$;\
$A \oplus B = \{1,5,6,7,11,12\}$.
\end{solution}

\begin{solution}{7.2}
The region consists of the two ``lune'' parts of the discs --- everything in
exactly one of $A, B$, excluding the central lens $A \cap B$ and the
exterior. This is the symmetric difference $A \oplus B$.
\end{solution}

\begin{solution}{7.3}
$|A \cup B| = 28 + 19 - 7 = 40$.
\end{solution}

\begin{solution}{7.4}
$x \in A \setminus (B \cup C)
\iff x \in A \wedge \neg(x \in B \vee x \in C)
\iff x \in A \wedge x \notin B \wedge x \notin C$
(De Morgan). Also
$x \in (A \setminus B) \cap (A \setminus C)
\iff (x \in A \wedge x \notin B) \wedge (x \in A \wedge x \notin C)$,
which by idempotence of $\wedge$ is the same condition. Hence the sets are
equal.
\end{solution}

\begin{solution}{7.5}
$A \cup (B \setminus A)
= A \cup (B \cap \overline{A})
= (A \cup B) \cap (A \cup \overline{A})
= (A \cup B) \cap U = A \cup B$,
using difference-as-intersection, distributivity, complement and identity
laws.
\end{solution}

\begin{solution}{7.6}
\begin{center}
\begin{tabular}{cc|ccc}
\toprule
$A$&$B$&$A\setminus B$&$A\setminus(A\setminus B)$&$A\cap B$\\
\midrule
1&1&0&1&1\\
1&0&1&0&0\\
0&1&0&0&0\\
0&0&0&0&0\\
\bottomrule
\end{tabular}
\end{center}
Columns 4 and 5 coincide.
\end{solution}

\begin{solution}{7.7}
If $A \subseteq B$: every element of $A \cup B$ is in $B$, so
$A \cup B \subseteq B \subseteq A \cup B$, giving equality; and
$A \subseteq A \cap B \subseteq A$ similarly. Conversely, if $A \cup B = B$
then $A \subseteq A \cup B = B$; if $A \cap B = A$ then
$A = A \cap B \subseteq B$.
\end{solution}

\begin{solution}{7.8}
Membership in a symmetric difference holds exactly when membership holds in
an \emph{odd} number of the operands: for two sets this is the definition.
Then $x \in (A \oplus B) \oplus C$ iff $x$ lies in an odd number of
$\{A,B,C\}$ --- check: it requires oddness between $(A \oplus B)$ and $C$,
and $x \in A \oplus B$ records oddness between $A$ and $B$; combining
parities gives overall parity. The same characterisation holds for
$A \oplus (B \oplus C)$, so the two sides are equal. (The eight-row
membership table gives the same conclusion mechanically.)
\end{solution}

\begin{solution}{7.9}
Take $X = A \oplus B$. Then, using associativity (Exercise 7.8),
$A \oplus (A \oplus B) = (A \oplus A) \oplus B = \varnothing \oplus B = B$.
Uniqueness: if $A \oplus X = A \oplus Y$, then XOR-ing both sides with $A$
gives $X = Y$. (The structure $(\powerset(U), \oplus)$ is a group in which
every element is its own inverse.)
\end{solution}

\begin{solution}{7.10}
Count the contribution of an element $x$ belonging to exactly $k \geq 1$ of
the three sets. Left side counts $x$ once. Right side counts it
$\binom{k}{1} - \binom{k}{2} + \binom{k}{3}$ times: for $k=1$: $1$; $k=2$:
$2 - 1 = 1$; $k=3$: $3 - 3 + 1 = 1$. Every element is counted once on each
side, so the totals agree. (Chapter~10 generalises this to $n$ sets.)
\end{solution}
